\chapter[GLMs and optimization]{Exponential families, GLMs, and constrained optimization}
\label{ch:glm-optimization}
The Poisson model in \path{poisson_spike_glm.m}
can be derived from a probabilistic model, rather than treated
as an unexplained numerical recipe.

\section{Poisson counts as an exponential family}
For count $Y\in\{0,1,\ldots\}$ and positive mean $\mu$,
\[
p(y\mid\mu)=\exp\bigl(y\log\mu-\mu-\log(y!)\bigr).
\]
Its natural parameter is $\eta=\log\mu$. For a row of
covariates $x_i^\top$, exposure duration $\Delta_i>0$,
and regression coefficients $\beta$, define
$\mu_i=\Delta_i\exp(x_i^\top\beta)$.
In matrix form with design $X$,
\[
\ell(\beta)=\sum_i[y_i(x_i^\top\beta+\log\Delta_i)
 -\Delta_i e^{x_i^\top\beta}-\log(y_i!)].
\]
\begin{proposition}[Concavity of the Poisson log likelihood]
For fixed nonnegative counts and strictly positive exposures,
\[
\nabla\ell=X^\top(y-\mu),\qquad
\nabla^2\ell=-X^\top W X,\quad W=\diag(\mu_i)\succ0.
\]
Thus $\ell$ is concave; if $X$ has full column rank, it is
strictly concave at every finite coefficient vector.
\end{proposition}
\begin{proof}
Differentiate the exponential using
\[\frac{\partial\mu_i}{\partial\beta}=\mu_i x_i.\]
For nonzero $v$, $-v^\top\nabla^2\ell v
=\sum_i\mu_i(x_i^\top v)^2$, positive under full rank.
\end{proof}
\emph{Strict concavity is not existence of a finite maximizer.}
For all-zero counts and an unpenalized intercept, likelihood
increases as the intercept tends to $-\infty$.

\section{IRLS as Fisher scoring}
Set $\eta=X\beta+\log\Delta$ and $\mu=e^\eta$.
Newton scoring solves, at the current iterate,
\[
(X^\top W X)\,\delta=X^\top(y-\mu),
\qquad \beta_{\mathrm{new}}=\beta+\delta.
\]
For known exposure offset $o=\log\Delta$, the full
working response is $z_{\mathrm{full}}=\eta+(y-\mu)/\mu$.
Weighted least squares on the design matrix $X$ must use
the offset-adjusted response:
\[
z_{\mathrm{WLS}}=z_{\mathrm{full}}-o
=X\beta+\frac{y-\mu}{\mu},\qquad
\beta_{\mathrm{new}}=(X^\top W X)^{-1}X^\top Wz_{\mathrm{WLS}},
\]
when the inverse exists. Alternatively fit $z_{\mathrm{full}}$
while holding $o$ fixed as an offset, not as a fitted predictor.
For the canonical Poisson family, the observed
negative Hessian agrees with expected Fisher information.
A quadratic ridge penalty adds $-\lambda P\beta$ to the score
and $\lambda P$ to the negative Hessian. If the algorithm
\emph{clips} $\eta$, or lacks line search, the numerical
iterations need not monotonically improve the original
unclipped likelihood. This is relevant to the present IRLS draft.

\section{Jensen, the ELBO, and EM monotonicity}
For latent variable $Z$ and any density $q(z)$ supported
where the joint model is positive, Jensen gives
\[
\log p_\theta(y)
=\log\int q(z)\frac{p_\theta(y,z)}{q(z)}\,dz
\ge \mathcal L(q,\theta)
:=\E_q[\log p_\theta(y,Z)]-\E_q[\log q(Z)].
\]
The identity
$\log p_\theta(y)=\mathcal L(q,\theta)
 +D_{\mathrm{KL}}(q\|p_\theta(z\mid y))$
shows exactly when the lower bound is tight.
Exact EM alternates an E-step setting $q$ to the posterior and
an M-step not decreasing the ELBO. This gives nondecreasing
observed-data likelihood, but \emph{not} a global optimum.
For phase--amplitude coupling, the Tort modulation index is
a normalized KL divergence between a binned phase-dependent
amplitude distribution and the uniform distribution; positivity
comes from Gibbs' inequality, while binning and surrogates
are additional statistical decisions.

\section{KKT and alternating least squares}
For $\min_\beta f(\beta)$ with differentiable inequalities
$g_j(\beta)\le0$ and affine equalities $A\beta=b$,
the KKT system consists of primal feasibility,
multipliers $\lambda_j\ge0$, complementary slackness
$\lambda_jg_j(\beta)=0$, and stationarity
\[
\nabla f(\beta)+\sum_j\lambda_j\nabla g_j(\beta)+A^\top\nu=0.
\]
\emph{Necessity} for a local minimizer requires a suitable
constraint qualification, such as LICQ.
\emph{Sufficiency} for global optimality holds when
$f$ and $g_j$ are convex differentiable functions and
the equality constraints are affine.
Slater's condition is a standard qualification under its
usual convex-program hypotheses. Lasso stationarity uses
subgradients at zero.
For a CP tensor decomposition, exact minimization over a factor
block cannot increase the unpenalized objective \emph{of that
specified block problem}; alternating such steps does not
make the multi-factor objective convex or guarantee global
optimality. The exploratory \path{tensor_cp_als.m}
includes ridge terms and does not track a monotonic objective;
a monotonicity claim for that code must be tested, not assumed.
% --- expanded edition (October 2026) ---
\section{Deviance, model comparison, and overdispersion}
The Poisson \emph{deviance} compares a fitted model with the saturated
model $\hat\mu_i=y_i$:
\[
D=2\sum_i\Bigl[y_i\log\frac{y_i}{\hat\mu_i}-(y_i-\hat\mu_i)\Bigr],
\]
with the convention $0\log0=0$. For nested models the deviance
difference is twice the log-likelihood ratio, approximately $\chi^2$
with degrees of freedom equal to the difference in parameters
(Result 52). Penalized comparisons use $\mathrm{AIC}=-2\ell+2p$ or
$\mathrm{BIC}=-2\ell+p\log n$; for spike trains, the effective $n$ in
BIC is ambiguous (bins? spikes? trials?), and cross-validated
log-likelihood is usually more defensible.

If the data are overdispersed relative to Poisson, the MLE of $\beta$
remains consistent when the mean model is correct, but its standard
errors from $(X^\top WX)^{-1}$ are too small. A quasi-Poisson correction
multiplies the covariance by the dispersion estimate
$\hat\phi=\frac{1}{n-p}\sum_i(y_i-\hat\mu_i)^2/\hat\mu_i$; the
sandwich estimator $(X^\top WX)^{-1}(X^\top\hat\Omega X)(X^\top WX)^{-1}$,
with $\hat\Omega=\diag((y_i-\hat\mu_i)^2)$, is robust to any variance
misspecification but not to temporal correlation of residuals, which
calls for block-based variants.

\section{Bernoulli GLMs and the logistic link}
In fine bins where at most one spike can occur, $y_i\in\{0,1\}$ and a
Bernoulli model with $\Prob(y_i=1)=\sigma(x_i^\top\beta)$,
$\sigma(u)=1/(1+e^{-u})$, is natural. The log-likelihood
$\sum_i[y_ix_i^\top\beta-\log(1+e^{x_i^\top\beta})]$ is concave with
Hessian $-X^\top\diag(p_i(1-p_i))X$. For small probabilities
$\sigma(u)\approx e^u$ and the model reduces to the Poisson GLM with
log link. The \emph{complementary log-log} link
$\Prob(y_i=1)=1-\exp(-\Delta e^{x_i^\top\beta})$ is exactly the probability
of at least one event in a bin for a Poisson process with rate
$e^{x_i^\top\beta}$, and is therefore the most faithful bin-level model
of an underlying continuous-time intensity.

\section{Sparse penalties and soft thresholding}
For orthonormal design ($X^\top X=I$) and squared loss, the lasso
problem $\min_\beta\frac12\|y-X\beta\|^2+\lambda\|\beta\|_1$ separates
into scalar problems with solution
\[
\hat\beta_j=\operatorname{sign}(z_j)\max(|z_j|-\lambda,0),\qquad z=X^\top y.
\]
\begin{proof}
For one coordinate minimize $\frac12(\beta-z)^2+\lambda|\beta|$. The
subgradient condition $0\in\beta-z+\lambda\,\partial|\beta|$ gives
$\beta=z-\lambda$ if $z>\lambda$, $\beta=z+\lambda$ if $z<-\lambda$, and
$\beta=0$ when $|z|\le\lambda$, since then $0\in-z+\lambda[-1,1]$.
\end{proof}
For general designs, coordinate descent cycles this update with partial
residuals; for GLMs it is applied inside each IRLS step. Lasso coupling
filters are attractive because they return exact zeros, but which
coupling coefficients are zero depends on correlations among regressors:
among two strongly correlated presynaptic neurons, the lasso tends to
select one arbitrarily. Zero coefficients are a property of the
estimator, not evidence of absent connections.

\section{An EM algorithm in closed form}
Suppose each trial's count is drawn from one of two Poisson rates,
$y_i\sim\pi\Poisson(\lambda_1)+(1-\pi)\Poisson(\lambda_2)$, for example
from an up and a down state. With responsibilities
\[
r_i=\frac{\pi\,p(y_i\mid\lambda_1)}{\pi\,p(y_i\mid\lambda_1)+(1-\pi)\,p(y_i\mid\lambda_2)}
\]
(E-step), the M-step maximizes the expected complete-data log-likelihood
and gives
\[
\pi=\frac1n\sum_ir_i,\qquad\lambda_1=\frac{\sum_ir_iy_i}{\sum_ir_i},\qquad
\lambda_2=\frac{\sum_i(1-r_i)y_i}{\sum_i(1-r_i)}.
\]
Each iteration does not decrease the observed-data likelihood (the ELBO
argument of this chapter), but the likelihood has a symmetric copy under
label swapping and may have degenerate maxima where one component
collapses onto a single observation. Several initializations and a
check that the fitted components are separated are required before
interpreting them as ``states''.

\section*{Exercises for this chapter}
\begin{exercise}
Show that the Poisson deviance is nonnegative and zero only for a
perfect fit.
\end{exercise}
\begin{exercise}
Derive the Hessian of the Bernoulli log-likelihood and show it is
negative semidefinite.
\end{exercise}
\begin{exercise}
Show that for small $\Delta e^{u}$ the complementary log-log
probability is approximately $\Delta e^u$.
\end{exercise}
\begin{exercise}
Derive the M-step formulas for the two-component Poisson mixture.
\end{exercise}
